% Current prompt marking throughout: x=(sr)^(-2), y=sxs,
% phi(x)=s, phi(y)=sr. This is not LM2's simultaneous marking.
\subsection{The normalized descent and a marking obstruction}

\paragraph{LC.1 (Proved): actual normalized descent frames at multiplier one.}
Put $\tau=sr$, $x=\tau^{-2}$, $y=sxs$, and let
$\varphi(x)=s$, $\varphi(y)=\tau$. For $\Phi=(1,f)\in\mathcal H$ set
\[
 C=\varphi(f),\qquad H=\varphi\bigl(f(z,x)\bigr).
\]
Then $C\in F$. Write $C^{\rm ab}=(p,q)$ in the basis $x,y$, put
$a=q-p$, and let $l\in\{0,1,2\}$ be defined by
$H\bmod F=r^l$. The descent in the prompt's stated marking has
the normalized pair and frames
\begin{align*}
 f_D&=y^a\theta(C)C^{-1}x^{-a},\qquad \theta(w)=sws^{-1},\\
 \beta_D&=x^aC,\qquad A_D=x^aCHr^{-l},\qquad \rho_D=-q.
\end{align*}
In particular these are actual pure frames, and the exponent $a$ is
uniquely specified; an unrecorded inner change is not being absorbed
into a signed coordinate.

\emph{Proof.}
Conjugation by $r$ cyclically permutes $x,y,z$. The normalized gentle
extension has $\Phi(r)=f^{-1}r$. Applying this twice, and using
$f(z,x)f(y,z)f(x,y)=1$, gives
\[
 \Phi(z)=f(z,x)z f(z,x)^{-1}.
\]
The kernel of $\varphi$ is the closed normal subgroup generated by
$x^2,z^3$. It is preserved by $\Phi$ and its inverse. Thus the descent
$\psi$ is an actual automorphism of $\widehat B$, with
\[
 \psi(s)=s,\qquad \psi(\tau)=C^{-1}\tau C,
 \qquad \psi(r)=HrH^{-1}.
\]
Since $f\in F'$, both $C,H$ lie in $\widehat B'$. Their images in
$S_3$ lie in $A_3=\langle r\rangle$. In particular $H$ centralizes
$r$ there, so $\psi$ induces the identity on $S_3$. The cusp equation
then says that the image of $C$ centralizes the transposition $\tau$.
Since $A_3\cap C_{S_3}(\tau)=1$, this proves $C\in F$.

The automorphism $\psi_0=\operatorname{Ad}_C\psi$ fixes $\tau$ and
has raw frames $C,CH$. Right multiplication of the latter by $r^{-l}$
makes it pure and does not change its conjugation of $r$. Now set
$\Phi_D=\operatorname{Ad}_{x^a}\psi_0$. This still fixes $\tau$.
Its pure $s$ frame is $x^aC$, with abelianization $(p+a,q)$.
The element
\[
 s\Phi_D(s)=y^a\theta(C)C^{-1}x^{-a}
\]
has abelianization $(q-p-a,p-q+a)$, so it belongs to $F'$ precisely
when $a=q-p$. The normalized representative then has
$\Phi_D(x)=x$, $\Phi_D(y)=f_D^{-1}yf_D$,
$\Phi_D(s)=sf_D$, and $\Phi_D(r)=f_D^{-1}r$.
These actual identities prove that it is the normalized gentle pair
representing the descended outer action. They also give the two
displayed frames and $\rho_D=-q$.

We record two integral consequences needed below. The image of
$C\in\widehat B'$ in $\widehat B^{\rm ab}=C_6$ gives
$p+q\equiv0\pmod3$. Further, $A_D\in F'$: the equality
$f_D=rA_Dr^{-1}A_D^{-1}$ implies
$(\operatorname{Ad}_r-1)A_D^{\rm ab}=0$. On
$F^{\rm ab}=\widehat{\mathbb Z}^2$ the determinant of this linear
operator is $3$. Its kernel is zero because multiplication by $3$
is injective on $\widehat{\mathbb Z}$.

\paragraph{LC.2 (Proved): the signed finite formula for $\Delta\circ D$.}
Use the matrices of R4,
\[
 U=\begin{pmatrix}1&1\\0&1\end{pmatrix},\quad
 S=\begin{pmatrix}0&1\\-1&0\end{pmatrix},\quad
 R_0=\begin{pmatrix}0&-1\\1&-1\end{pmatrix}.
\]
Let
\[
 \mathsf C=f(S,U^{-1}),\qquad
 \mathsf H=f(R_0^{-1},S)
       \quad\hbox{in }\mathrm{SL}_2(\mathbb Z/12).
\]
These are the commutator lifts of $C,H$; they must not be substituted
for the signed pure lifts without a sign correction. With $a,l$ as
in LC.1,
\[
 a\equiv l\pmod3,\qquad
 q_{12}(\beta_D)=(-I)^aU^{2a}\mathsf C,\qquad
 q_{12}(A_D)=(-I)^aU^{2a}\mathsf C\mathsf H R_0^{-l}.
\]
Modulo three write
$\mathsf C=S^j$ and $\mathsf C\mathsf H=(-I)^e$; such
$j\in\mathbb Z/4$ and $e\in\mathbb Z/2$ exist for every actual
gentle pair. Then
\[
 \boxed{\ \Delta(D\Phi)=a+e+j\pmod2.\ }
\]
The parity of $a$ is itself finite: modulo four,
$\mathsf C=(-U^2)^a$. Thus the displayed formula is an exact
finite evaluation formula on actual gentle pairs.

\emph{Proof.}
The central extension $\widehat{\mathrm{SL}_2(\mathbb Z)}\to
\widehat B$ has kernel $\{\pm I\}$. Its abelianization is $C_{12}$,
in which $U,L$ have class $1$, $S=ULU$ has class $3$,
$-I$ has class $6$, and $R_0=-UL$ has class $8$.
Its derived subgroup maps isomorphically onto $\widehat B'$:
the two possible lifts of an element of $\widehat B'$ have
abelianizations $0,6$, and exactly one is in the derived subgroup.
The words defining $\mathsf C,\mathsf H$ choose this latter lift.
For $H$, the lift of $\varphi(z)$ differs from $R_0^{-1}$ by $-I$,
which does not change evaluation of the commutator word $f$.

The signed pure lift of $C$ instead has abelianization $2(p+q)$.
Since $p+q\equiv0\pmod3$ and
$p+q\equiv q-p=a\pmod2$, it differs from the commutator lift by
$(-I)^a$. This proves the formula for $\beta_D$. The word
$U^{2a}\mathsf C\mathsf H R_0^{-l}$ has abelianization
$2a-8l$, while the signed lift of $A_D\in F'$ has abelianization
zero. The two lifts differ by a sign. Consequently
$a\equiv l\pmod3$, and that sign is again $(-I)^a$, since
$(a-4l)/3\equiv a\pmod2$. This proves the second formula.

Inserting $\rho_D=-q$ into the R4 cusp correction gives, modulo
three,
\[
 W_D=(-I)^a\mathsf C^{-1},\qquad
 T_D=(-I)^aR_0^l(\mathsf C\mathsf H)^{-1}.
\]
Indeed the remaining upper-unipotent exponent is
$-2a-8q=2p-10q\equiv2(p+q)=0\pmod3$.
The R4 centralizer condition implies $\mathsf C\in\langle S\rangle$.
Both $\mathsf C,\mathsf H$ modulo three lie in
$\mathrm{SL}_2(\mathbb F_3)'=Q_8$. Since
$C_{\mathrm{SL}_2(\mathbb F_3)}(R_0)\cap Q_8=\{\pm I\}$,
their product is $(-I)^e$. The sign character of $T_D$ is $a+e$,
while that of $W_D$ is $j$, proving the boxed formula. Finally
$q_4(\beta_D)=I$ in R4 and $U^{-2}=U^2$ modulo four give
$\mathsf C=(-U^2)^a$ modulo four.

\paragraph{LC.3 (Refuted): the claim that the stated $D$ fixes the Galois image.}
For the precise current marking in LC.1 one has
\[
 \boxed{\ D(\sigma)=c\sigma c
       \quad(\sigma\in G_{\mathbb Q^{\rm ab}}),\qquad
       \rho(D\sigma)=-\rho(\sigma).\ }
\]
There exist such $\sigma$ with $D(\sigma)\ne\sigma$.
Consequently the literal inclusion
$\widehat{\mathrm{GT}}\subseteq\operatorname{Fix}(D)$ is false
with this marking. Identifying this fixed subgroup with the manuscript's
arithmetically compatible subgroup without a marking change is also
false.

\emph{Proof.}
Let $u=\tau^{-1}$ and $s=uvu$. Write $g=\beta_\sigma^{-1}$.
With positive real radical roots, Nakamura--Schneps prove
\[
 g^{\rm ab}=(\rho_2,\rho_2),\qquad \rho=\rho_2
\]
\cite[Proposition~2.4 and Definition~5.1]{LiftNakamuraSchneps2000}.
The two actual braid identities in the authors' opened report are
\begin{align*}
 g(u^2,v^2)&=f(u^2,s)u^{4\rho_2},\\
 f(u,s)&=s^{\rho_3-2\rho_2}f(u^2,s)u^{6\rho_2-3\rho_3}
\end{align*}
\cite[Theorem~2, $(\mathrm{GF}_1)$; Theorem~1, $(\mathrm V)$]
{LiftTsunogaiNakamura2002}.
For $\sigma\in G_{\mathbb Q^{\rm ab}}$ the Kummer exponents
$\rho_2,\rho_3$ are even: $\sqrt2,\sqrt3\in\mathbb Q^{\rm ab}$.
Passing to $\widehat B$, where $s^2=1$, therefore gives
\[
 f(u,s)=g\,u^{2\rho_2-3\rho_3}.
\]
Let $c_B$ denote the actual complex-conjugation automorphism of
$\widehat B$: $c_B(u)=u^{-1}$, $c_B(v)=v^{-1}$, $c_B(s)=s$.
The two-cycle relation and the preceding equality yield
\begin{align*}
 C=f(s,u^{-1})
   &=c_B(f(u,s))^{-1}\\
   &=u^{2\rho_2-3\rho_3}c_B(\beta_\sigma)
     =\tau^{-2\rho_2+3\rho_3}c_B(\beta_\sigma).
\end{align*}
Since $c_B$ negates $F^{\rm ab}$ and $\beta_\sigma^{\rm ab}
=(-\rho_2,-\rho_2)$, it follows that
\[
 C^{\rm ab}=(2\rho_2-3\rho_3/2,\rho_2),\qquad
 a=-\rho_2+3\rho_3/2.
\]
The correction of LC.1 now cancels exactly:
\[
 \beta_D=x^aC=c_B(\beta_\sigma).
\]
Both $\Phi_D$ and $c_B\Phi_\sigma c_B$ fix $\tau$ literally;
the latter has the pure $s$ frame $c_B(\beta_\sigma)$.
They agree on $s,\tau$, which generate $\widehat B$, and hence are
equal as actual automorphisms. This proves the boxed identities.

For an actual counterexample put $K=\mathbb Q(\zeta_8)$ and
$L=K(\alpha)$ with $\alpha=2^{1/8}>0$. The field
$\mathbb Q(2^{1/4})$ is not Galois over $\mathbb Q$:
$T^4-2$ is Eisenstein and its real root field does not contain
the nonreal root $i2^{1/4}$. It therefore is not contained in
the abelian extension $K/\mathbb Q$. Since $\mu_4\subset K$
and $\alpha^4=\sqrt2\in K$, the extension $K(\alpha)/K$
is cyclic of degree dividing four; degree at most two would
put $\alpha^2=2^{1/4}$ in $K$. Thus $[L:K]=4$ and
$[L:\mathbb Q]=16$. It is the splitting field of $T^8-2$.

Write its automorphisms as $(v,b)$, meaning
$\zeta_8\mapsto\zeta_8^v$ and
$\alpha\mapsto\zeta_8^b\alpha$. Their exact set and law are
\[
 \begin{gathered}
 v\in\{1,3,5,7\},\quad b\in\mathbb Z/8,\qquad
 b\equiv\begin{cases}0&v=1,7,\\1&v=3,5\end{cases}\pmod2,\\
 (v,b)(v',b')=(vv',b+vb').
 \end{gathered}
\]
The parity condition is exactly the image of
$\alpha^4=\sqrt2=\zeta_8+\zeta_8^{-1}$. Exact enumeration
gives the commutator subgroup $\{(1,0),(1,2),(1,4),(1,6)\}$.
Choose an actual $\sigma\in G_{\mathbb Q^{\rm ab}}$ restricting
to $(1,2)$; this is possible since restriction maps the closed
commutator subgroup of $G_{\mathbb Q}$ onto the commutator
subgroup of every finite Galois quotient. Then
$\rho_2(\sigma)\equiv2\pmod8$, whereas
$\rho_2(D\sigma)\equiv6\pmod8$. Complex conjugation is
$(7,0)$, and conjugation by it indeed sends $(1,2)$ to $(1,6)$.
This proves $D\sigma\ne\sigma$ in an actual degree-sixteen
Galois quotient. All group laws, the commutator subgroup, an
explicit commutator for $(1,2)$, and the distinction modulo eight
are checked in \texttt{check\_descent\_marking.py}.

\paragraph{LC.4 (Proved): the exact conversion from the LM2 marking.}
The old LM2 generators are
\[
 x_L=(rs)^2=y^{-1},\qquad y_L=(sr)^2=x^{-1}.
\]
Consequently its map $\varphi_L(x_L)=s$, $\varphi_L(y_L)=\tau$,
when written on the current generators, is
\[
 \varphi_L(x)=\tau^{-1},\qquad \varphi_L(y)=s.
\]
Put $k(x)=y^{-1}$, $k(y)=x^{-1}$. Then
$\varphi_L=\varphi\circ k$ and $k=\theta\circ c$ on $F$.
Every gentle outer class commutes with the geometric interchange
$\theta$. Writing $D_L$ for the old descent followed by restriction
to the same embedded pure subgroup therefore gives
\[
 D_L(\Phi)=D(c\Phi c)=cD(\Phi)c.
\]
Here $D(c)=c$ follows directly from
$\psi(s)=s$, $\psi(\tau)=\tau^{-1}$. Thus the corrected
endomorphism for the manuscript fixed-point convention is
$D^*=\operatorname{Ad}_c\circ D$. The two functions
$\Delta\circ D^*$ and $\Delta\circ D$ agree by R4's
conjugation invariance, but their fixed subgroups are different.
This change is required before transporting equations (30), NT
compatibilities, or the assertion that the Galois image is fixed.

\emph{Proof.}
The displayed word substitutions prove the marking formulas
without a choice of path. Descent along $\varphi\circ k$ is descent
along $\varphi$ after conjugating the source automorphism by $k$.
The geometric interchange commutes with gentle outer classes, so
this conjugation is the same as conjugation by $c$ at the outer
level. Descent and restriction are homomorphisms on these outer
groups. Together with $D(c)=c$ this gives the equality of
endomorphisms. The final equality of characters is precisely the
conjugation invariance already proved in R4. LC.3 independently
checks the resulting nontrivial twist on the actual Galois action,
including all cusp and frame normalizations.

\paragraph{LC.5 (Boundary of the calculation, not a profinite conclusion).}
The exact finite formula LC.2 makes a direct test possible. The
six substitutions of the same $f$ into the pairs
\[
 (S,U^{-1}),\ (U^{-1},Z),\ (Z,S),\
 (U^{-1},S),\ (Z,U^{-1}),\ (S,Z),\qquad Z=(SU^{-1})^{-1},
\]
give a finite group of order $8153726976$, with derived subgroup of
order $56623104$. Exact membership, signed-frame conditions, and
the two-cycle and three-cycle equations retain twelve candidates
of each value of the right-hand side of LC.2. Including the original
inclusion evaluation $(U^2,L^2)$ and its required R3 value for a fixed
point still does not force zero. The opposite source sign required by
the corrected locus $D(\Phi)=c\Phi c$ retains the same counts.
These computations do not establish
that any retained candidate has a compatible profinite lift. Their
only purpose is to identify why this finite-evaluation attempt has
not proved $\Delta\circ D=0$ or the asserted fixed-subgroup inclusion.
The actual remaining question requires further profinite or pentagon
constraints. The exact implementation and its output are retained
as \texttt{check\_descent\_evaluations.py} and its records.
